% !TEX root = ../report.tex
%%
%% Examples of frequently used elements.
%% This chapter is only a cheat sheet: copy what you need into your own
%% chapters and then delete the line % !TEX root = ../report.tex
%%
%% Examples of frequently used elements.
%% This chapter is only a cheat sheet: copy what you need into your own
%% chapters and then delete the line % !TEX root = ../report.tex
%%
%% Examples of frequently used elements.
%% This chapter is only a cheat sheet: copy what you need into your own
%% chapters and then delete the line % !TEX root = ../report.tex
%%
%% Examples of frequently used elements.
%% This chapter is only a cheat sheet: copy what you need into your own
%% chapters and then delete the line \input{chapters/99_beispiele} in
%% report.tex.
%%
\chapter{\DeEn{Beispiele (bitte löschen)}{Examples (please delete)}}\label{chap:beispiele}

\section{Zitate und Verweise}
Ein Literaturverweis: Grundlagen des statistischen Lernens finden sich in
\cite{hastie2009}, das Vorgehensmodell CRISP-DM in \cite{chapman2000}.
Mit natbib geht auch \citet{pedregosa2011} (Autor im Text) oder ein Verweis
mit Seitenangabe \cite[S.~42]{hastie2009}.
Querverweise: \autoref{fig:lernkurve}, \autoref{tab:metriken} und
\autoref{lst:sklearn}.

\section{Abkürzungen, Glossar und Symbole}
Beim ersten Auftreten wird eine Abkürzung ausgeschrieben: \gls{ac:ml}.
Danach nur noch die Kurzform: \gls{ac:ml}. Ein Glossarbegriff: \gls{gl:feature}.
Ein Symbol: Die Verlustfunktion \gls{sy:loss} wird minimiert.

\section{Abbildungen}
\autoref{fig:lernkurve} zeigt ein mit \texttt{pgfplots} direkt in \LaTeX{}
erzeugtes Diagramm. Bilddateien (PDF, PNG, JPG) legt man in \texttt{figures/}
ab und bindet sie z.\,B. so ein:
\begin{lstlisting}[language={[LaTeX]TeX}, numbers=none]
\includenamedimage[0.7]{dateiname}{Bildunterschrift}
\end{lstlisting}

\begin{figure}[htbp]
  \centering
  \begin{tikzpicture}
    \begin{axis}[
        width=0.8\textwidth, height=6cm,
        xlabel={\DeEn{Epoche}{Epoch}}, ylabel={\DeEn{Verlust}{Loss}},
        grid=major, legend pos=north east]
      \addplot[thick, color=IAccent, domain=1:50, samples=50] {exp(-x/10) + 0.1};
      \addlegendentry{Training}
      \addplot[thick, dashed, color=gray, domain=1:50, samples=50] {exp(-x/12) + 0.15 + 0.002*x};
      \addlegendentry{\DeEn{Validierung}{Validation}}
    \end{axis}
  \end{tikzpicture}
  \caption{\DeEn{Beispielhafte Lernkurve}{Example learning curve}}
  \label{fig:lernkurve}
\end{figure}

\section{Tabellen}
\begin{table}[htbp]
  \centering
  \caption{\DeEn{Vergleich von Modellen}{Comparison of models}}
  \label{tab:metriken}
  \begin{tabular}{lSS}
    \toprule
    {Modell} & {Accuracy} & {F1-Score} \\
    \midrule
    Logistische Regression & 0.842 & 0.811 \\
    Random Forest          & 0.887 & 0.864 \\
    Gradient Boosting      & 0.901 & 0.879 \\
    \bottomrule
  \end{tabular}
\end{table}

\section{Quellcode}
Kurzer Code im Text: \lstinline|model.fit(X, y)|.

\begin{lstlisting}[caption={\DeEn{Training eines Modells mit scikit-learn}{Training a model with scikit-learn}}, label={lst:sklearn}]
from sklearn.ensemble import RandomForestClassifier
from sklearn.model_selection import train_test_split

# Daten aufteilen und Modell trainieren
X_train, X_test, y_train, y_test = train_test_split(X, y, test_size=0.2, random_state=42)
model = RandomForestClassifier(n_estimators=100, random_state=42)
model.fit(X_train, y_train)
print(f"Accuracy: {model.score(X_test, y_test):.3f}")
\end{lstlisting}

\section{Formeln und Einheiten}
Die mittlere quadratische Abweichung ist definiert als
\begin{equation}
  \mathrm{MSE} = \frac{1}{n} \sum_{i=1}^{n} \left( y_i - \hat{y}_i \right)^2 .
  \label{eq:mse}
\end{equation}
Einheiten mit \texttt{siunitx}: Das Training dauerte \SI{12.5}{\minute} auf
\num{1250000} Datensätzen.

\section{Notizen}
\todo{Nur im Entwurfsmodus sichtbar.}
Mit \verb|\todo{...}| lassen sich Anmerkungen am Rand setzen.
 in
%% report.tex.
%%
\chapter{\DeEn{Beispiele (bitte löschen)}{Examples (please delete)}}\label{chap:beispiele}

\section{Zitate und Verweise}
Ein Literaturverweis: Grundlagen des statistischen Lernens finden sich in
\cite{hastie2009}, das Vorgehensmodell CRISP-DM in \cite{chapman2000}.
Mit natbib geht auch \citet{pedregosa2011} (Autor im Text) oder ein Verweis
mit Seitenangabe \cite[S.~42]{hastie2009}.
Querverweise: \autoref{fig:lernkurve}, \autoref{tab:metriken} und
\autoref{lst:sklearn}.

\section{Abkürzungen, Glossar und Symbole}
Beim ersten Auftreten wird eine Abkürzung ausgeschrieben: \gls{ac:ml}.
Danach nur noch die Kurzform: \gls{ac:ml}. Ein Glossarbegriff: \gls{gl:feature}.
Ein Symbol: Die Verlustfunktion \gls{sy:loss} wird minimiert.

\section{Abbildungen}
\autoref{fig:lernkurve} zeigt ein mit \texttt{pgfplots} direkt in \LaTeX{}
erzeugtes Diagramm. Bilddateien (PDF, PNG, JPG) legt man in \texttt{figures/}
ab und bindet sie z.\,B. so ein:
\begin{lstlisting}[language={[LaTeX]TeX}, numbers=none]
\includenamedimage[0.7]{dateiname}{Bildunterschrift}
\end{lstlisting}

\begin{figure}[htbp]
  \centering
  \begin{tikzpicture}
    \begin{axis}[
        width=0.8\textwidth, height=6cm,
        xlabel={\DeEn{Epoche}{Epoch}}, ylabel={\DeEn{Verlust}{Loss}},
        grid=major, legend pos=north east]
      \addplot[thick, color=IAccent, domain=1:50, samples=50] {exp(-x/10) + 0.1};
      \addlegendentry{Training}
      \addplot[thick, dashed, color=gray, domain=1:50, samples=50] {exp(-x/12) + 0.15 + 0.002*x};
      \addlegendentry{\DeEn{Validierung}{Validation}}
    \end{axis}
  \end{tikzpicture}
  \caption{\DeEn{Beispielhafte Lernkurve}{Example learning curve}}
  \label{fig:lernkurve}
\end{figure}

\section{Tabellen}
\begin{table}[htbp]
  \centering
  \caption{\DeEn{Vergleich von Modellen}{Comparison of models}}
  \label{tab:metriken}
  \begin{tabular}{lSS}
    \toprule
    {Modell} & {Accuracy} & {F1-Score} \\
    \midrule
    Logistische Regression & 0.842 & 0.811 \\
    Random Forest          & 0.887 & 0.864 \\
    Gradient Boosting      & 0.901 & 0.879 \\
    \bottomrule
  \end{tabular}
\end{table}

\section{Quellcode}
Kurzer Code im Text: \lstinline|model.fit(X, y)|.

\begin{lstlisting}[caption={\DeEn{Training eines Modells mit scikit-learn}{Training a model with scikit-learn}}, label={lst:sklearn}]
from sklearn.ensemble import RandomForestClassifier
from sklearn.model_selection import train_test_split

# Daten aufteilen und Modell trainieren
X_train, X_test, y_train, y_test = train_test_split(X, y, test_size=0.2, random_state=42)
model = RandomForestClassifier(n_estimators=100, random_state=42)
model.fit(X_train, y_train)
print(f"Accuracy: {model.score(X_test, y_test):.3f}")
\end{lstlisting}

\section{Formeln und Einheiten}
Die mittlere quadratische Abweichung ist definiert als
\begin{equation}
  \mathrm{MSE} = \frac{1}{n} \sum_{i=1}^{n} \left( y_i - \hat{y}_i \right)^2 .
  \label{eq:mse}
\end{equation}
Einheiten mit \texttt{siunitx}: Das Training dauerte \SI{12.5}{\minute} auf
\num{1250000} Datensätzen.

\section{Notizen}
\todo{Nur im Entwurfsmodus sichtbar.}
Mit \verb|\todo{...}| lassen sich Anmerkungen am Rand setzen.
 in
%% report.tex.
%%
\chapter{\DeEn{Beispiele (bitte löschen)}{Examples (please delete)}}\label{chap:beispiele}

\section{Zitate und Verweise}
Ein Literaturverweis: Grundlagen des statistischen Lernens finden sich in
\cite{hastie2009}, das Vorgehensmodell CRISP-DM in \cite{chapman2000}.
Mit natbib geht auch \citet{pedregosa2011} (Autor im Text) oder ein Verweis
mit Seitenangabe \cite[S.~42]{hastie2009}.
Querverweise: \autoref{fig:lernkurve}, \autoref{tab:metriken} und
\autoref{lst:sklearn}.

\section{Abkürzungen, Glossar und Symbole}
Beim ersten Auftreten wird eine Abkürzung ausgeschrieben: \gls{ac:ml}.
Danach nur noch die Kurzform: \gls{ac:ml}. Ein Glossarbegriff: \gls{gl:feature}.
Ein Symbol: Die Verlustfunktion \gls{sy:loss} wird minimiert.

\section{Abbildungen}
\autoref{fig:lernkurve} zeigt ein mit \texttt{pgfplots} direkt in \LaTeX{}
erzeugtes Diagramm. Bilddateien (PDF, PNG, JPG) legt man in \texttt{figures/}
ab und bindet sie z.\,B. so ein:
\begin{lstlisting}[language={[LaTeX]TeX}, numbers=none]
\includenamedimage[0.7]{dateiname}{Bildunterschrift}
\end{lstlisting}

\begin{figure}[htbp]
  \centering
  \begin{tikzpicture}
    \begin{axis}[
        width=0.8\textwidth, height=6cm,
        xlabel={\DeEn{Epoche}{Epoch}}, ylabel={\DeEn{Verlust}{Loss}},
        grid=major, legend pos=north east]
      \addplot[thick, color=IAccent, domain=1:50, samples=50] {exp(-x/10) + 0.1};
      \addlegendentry{Training}
      \addplot[thick, dashed, color=gray, domain=1:50, samples=50] {exp(-x/12) + 0.15 + 0.002*x};
      \addlegendentry{\DeEn{Validierung}{Validation}}
    \end{axis}
  \end{tikzpicture}
  \caption{\DeEn{Beispielhafte Lernkurve}{Example learning curve}}
  \label{fig:lernkurve}
\end{figure}

\section{Tabellen}
\begin{table}[htbp]
  \centering
  \caption{\DeEn{Vergleich von Modellen}{Comparison of models}}
  \label{tab:metriken}
  \begin{tabular}{lSS}
    \toprule
    {Modell} & {Accuracy} & {F1-Score} \\
    \midrule
    Logistische Regression & 0.842 & 0.811 \\
    Random Forest          & 0.887 & 0.864 \\
    Gradient Boosting      & 0.901 & 0.879 \\
    \bottomrule
  \end{tabular}
\end{table}

\section{Quellcode}
Kurzer Code im Text: \lstinline|model.fit(X, y)|.

\begin{lstlisting}[caption={\DeEn{Training eines Modells mit scikit-learn}{Training a model with scikit-learn}}, label={lst:sklearn}]
from sklearn.ensemble import RandomForestClassifier
from sklearn.model_selection import train_test_split

# Daten aufteilen und Modell trainieren
X_train, X_test, y_train, y_test = train_test_split(X, y, test_size=0.2, random_state=42)
model = RandomForestClassifier(n_estimators=100, random_state=42)
model.fit(X_train, y_train)
print(f"Accuracy: {model.score(X_test, y_test):.3f}")
\end{lstlisting}

\section{Formeln und Einheiten}
Die mittlere quadratische Abweichung ist definiert als
\begin{equation}
  \mathrm{MSE} = \frac{1}{n} \sum_{i=1}^{n} \left( y_i - \hat{y}_i \right)^2 .
  \label{eq:mse}
\end{equation}
Einheiten mit \texttt{siunitx}: Das Training dauerte \SI{12.5}{\minute} auf
\num{1250000} Datensätzen.

\section{Notizen}
\todo{Nur im Entwurfsmodus sichtbar.}
Mit \verb|\todo{...}| lassen sich Anmerkungen am Rand setzen.
 in
%% report.tex.
%%
\chapter{\DeEn{Beispiele (bitte löschen)}{Examples (please delete)}}\label{chap:beispiele}

\section{Zitate und Verweise}
Ein Literaturverweis: Grundlagen des statistischen Lernens finden sich in
\cite{hastie2009}, das Vorgehensmodell CRISP-DM in \cite{chapman2000}.
Mit natbib geht auch \citet{pedregosa2011} (Autor im Text) oder ein Verweis
mit Seitenangabe \cite[S.~42]{hastie2009}.
Querverweise: \autoref{fig:lernkurve}, \autoref{tab:metriken} und
\autoref{lst:sklearn}.

\section{Abkürzungen, Glossar und Symbole}
Beim ersten Auftreten wird eine Abkürzung ausgeschrieben: \gls{ac:ml}.
Danach nur noch die Kurzform: \gls{ac:ml}. Ein Glossarbegriff: \gls{gl:feature}.
Ein Symbol: Die Verlustfunktion \gls{sy:loss} wird minimiert.

\section{Abbildungen}
\autoref{fig:lernkurve} zeigt ein mit \texttt{pgfplots} direkt in \LaTeX{}
erzeugtes Diagramm. Bilddateien (PDF, PNG, JPG) legt man in \texttt{figures/}
ab und bindet sie z.\,B. so ein:
\begin{lstlisting}[language={[LaTeX]TeX}, numbers=none]
\includenamedimage[0.7]{dateiname}{Bildunterschrift}
\end{lstlisting}

\begin{figure}[htbp]
  \centering
  \begin{tikzpicture}
    \begin{axis}[
        width=0.8\textwidth, height=6cm,
        xlabel={\DeEn{Epoche}{Epoch}}, ylabel={\DeEn{Verlust}{Loss}},
        grid=major, legend pos=north east]
      \addplot[thick, color=IAccent, domain=1:50, samples=50] {exp(-x/10) + 0.1};
      \addlegendentry{Training}
      \addplot[thick, dashed, color=gray, domain=1:50, samples=50] {exp(-x/12) + 0.15 + 0.002*x};
      \addlegendentry{\DeEn{Validierung}{Validation}}
    \end{axis}
  \end{tikzpicture}
  \caption{\DeEn{Beispielhafte Lernkurve}{Example learning curve}}
  \label{fig:lernkurve}
\end{figure}

\section{Tabellen}
\begin{table}[htbp]
  \centering
  \caption{\DeEn{Vergleich von Modellen}{Comparison of models}}
  \label{tab:metriken}
  \begin{tabular}{lSS}
    \toprule
    {Modell} & {Accuracy} & {F1-Score} \\
    \midrule
    Logistische Regression & 0.842 & 0.811 \\
    Random Forest          & 0.887 & 0.864 \\
    Gradient Boosting      & 0.901 & 0.879 \\
    \bottomrule
  \end{tabular}
\end{table}

\section{Quellcode}
Kurzer Code im Text: \lstinline|model.fit(X, y)|.

\begin{lstlisting}[caption={\DeEn{Training eines Modells mit scikit-learn}{Training a model with scikit-learn}}, label={lst:sklearn}]
from sklearn.ensemble import RandomForestClassifier
from sklearn.model_selection import train_test_split

# Daten aufteilen und Modell trainieren
X_train, X_test, y_train, y_test = train_test_split(X, y, test_size=0.2, random_state=42)
model = RandomForestClassifier(n_estimators=100, random_state=42)
model.fit(X_train, y_train)
print(f"Accuracy: {model.score(X_test, y_test):.3f}")
\end{lstlisting}

\section{Formeln und Einheiten}
Die mittlere quadratische Abweichung ist definiert als
\begin{equation}
  \mathrm{MSE} = \frac{1}{n} \sum_{i=1}^{n} \left( y_i - \hat{y}_i \right)^2 .
  \label{eq:mse}
\end{equation}
Einheiten mit \texttt{siunitx}: Das Training dauerte \SI{12.5}{\minute} auf
\num{1250000} Datensätzen.

\section{Notizen}
\todo{Nur im Entwurfsmodus sichtbar.}
Mit \verb|\todo{...}| lassen sich Anmerkungen am Rand setzen.
